% proofs.tex — Supplementary Material: Proofs (condensed version for review)
% Included in main.tex via \ifincludeproofs; also compiled standalone via supplement.tex.
% Theorem numbers refer to the shared counter of main.tex.

%==============================================================================
\section{Tool Inventory}\label{sec:S0}

\textbf{Tool S.1 (Rank inequality; \cite[Lemma~2.2]{bai2010spectral}).} For Hermitian $d\times d$ matrices $A,B$, $\|F^A-F^B\|_\infty\le\operatorname{rank}(A-B)/d$.

\textbf{Tool S.2 (Trace convexity chain; \cite{davis1957schwarz,fan1949theorem,karamata1932inegalite,marshall1979inequalities,bhatia1997matrix}).} Ky Fan majorization; Karamata; Davis: $\varphi$ convex $\Rightarrow$ $M\mapsto\operatorname{tr}\varphi(M)$ convex.

\textbf{Tool S.3 (Extreme eigenvalues of sample covariance matrices; \cite{bai1993limit}).} Let $Z\in\mathbb R^{d\times M}$ have i.i.d.\ entries drawn from a fixed distribution with mean $0$, variance $1$, finite fourth moment, and $d/M\to y\in(0,\infty)$. Then the extreme eigenvalues of $\frac1MZZ^\top$ converge in probability to the edges of $\mathrm{MP}(y)$: $\lambda_{\max}\to(1+\sqrt y)^2$; for $y<1$, $\lambda_{\min}\to(1-\sqrt y)^2$. For $y>1$, the Gram form gives directly: $\lambda_{\min}\big(\frac1MZ^\top Z\big)\to(\sqrt y-1)^2>0$ in probability—implying simultaneously that $Z$ has full column rank with probability tending to one and that the smallest positive eigenvalue of the sample covariance converges.

%==============================================================================
\section{Proof of Lemma~\ref{main-lem:unbiased}}\label{sec:S1}

$\mathbb E[S_{r,N}\mid X_N,w]=\frac1N\sum_i\mathbb E[n_i\mid X_N,w]\,x_ix_i^\top=\sum_iw_ix_ix_i^\top=S_{w,N}$, by linearity only. $\square$

%==============================================================================
\section{Proof of Theorem~\ref{main-thm:jensen} }\label{sec:S2}

The count vector of any scheme takes finitely many values $\boldsymbol n^{(1)},\dots,\boldsymbol n^{(L)}$ with probabilities $p_\ell$. By Lemma~\ref{main-lem:unbiased}, $\sum_\ell p_\ell S(\boldsymbol n^{(\ell)})=S_{w,N}$, and ordinary finite Jensen gives
\[
\mathbb E\big[F(S_{r,N})\,\big|\,X_N,w\big]=\sum_{\ell}p_\ell F\big(S(\boldsymbol n^{(\ell)})\big)\ \ge\ F(S_{w,N}).
\]
For strictly convex $F$, equality forces $S(\boldsymbol n^{(\ell)})$ identical across the support, i.e.\ $S_{r,N}=S_{w,N}$ a.s.

\textbf{Corollary} (a) Tool S.2(the Davis trace-convexity step); for $q>2$, $x\mapsto x^q$ is not operator convex on $(0,\infty)$, whereas trace convexity still applies. (b) $\lambda_{\max}(M)=\sup_{\|v\|=1}v^\top Mv$, a supremum of linear functions. (c) $-\log(x+\varepsilon)$ is convex on $[0,\infty)$; the conclusion reads $\mathbb E\log\det(S_{r,N}+\varepsilon I)\le\log\det(S_{w,N}+\varepsilon I)$, with the convention $0\log0=0$ throughout. $\square$

%==============================================================================
\section{Proof of Proposition~\ref{main-prop:cost} and Proposition~\ref{main-prop:popcost}}\label{sec:S3}

\textbf{Proposition~\ref{main-prop:cost}.} With $\delta_i=n_i/N-w_i$, $\Delta:=S_{r,N}-S_{w,N}=\sum_i\delta_ix_ix_i^\top$:
\[
\operatorname{tr}S_{r,N}^2-\operatorname{tr}S_{w,N}^2=2\operatorname{tr}(S_{w,N}\Delta)+\operatorname{tr}\Delta^2,\qquad \operatorname{tr}\Delta^2=\|\Delta\|_F^2.
\]
The linear term vanishes in conditional mean (Lemma~\ref{main-lem:unbiased}), so
\[
\mathbb E\big[\operatorname{tr}S_{r,N}^2-\operatorname{tr}S_{w,N}^2\,\big|\,X_N,w\big]=\mathbb E\big[\|\Delta\|_F^2\,\big|\,X_N,w\big]=\sum_{i,j}\mathbb E[\delta_i\delta_j]G_{ij}^2=\frac1{N^2}\operatorname{tr}\big[(G\circ G)\,\operatorname{Cov}(\boldsymbol n\mid X_N,w)\big],
\]
nonnegative since $G\circ G$ and the covariance are PSD (Schur product theorem).

\textbf{Multinomial specialization.} With $\mathbb E[n_i^2]=Nw_i(1-w_i)+N^2w_i^2$ and $\mathbb E[n_in_j]=N(N-1)w_iw_j$ ($i\neq j$), substitution gives $\frac1N[\sum_iw_i\|x_i\|^4-\operatorname{tr}S_{w,N}^2]$. For the equivalent geometric form, set $Y_i:=x_ix_i^\top$: since $\sum_iw_i=1$ and $\sum_iw_iY_i=S_{w,N}$,
\[
\sum_iw_i\big\|Y_i-S_{w,N}\big\|_F^2=\sum_iw_i\|Y_i\|_F^2-\Big\|\sum_iw_iY_i\Big\|_F^2,
\]
so the cost equals $\frac1N\sum_iw_i\|x_ix_i^\top-S_{w,N}\|_F^2$, which vanishes iff $x_ix_i^\top=S_{w,N}$ for every $i$ with $w_i>0$ (covering all-zero particles and the $\pm v$ collinear case; in the latter, $S_{r,N}=\frac1N(\sum_in_i)vv^\top=S_{w,N}$ a.s.\ even with random counts). $\square$

\textbf{Proposition~\ref{main-prop:popcost}.} For $x_i=\Sigma^{1/2}z_i$ with kurtosis $\kappa_4$:
\[
\mathbb E\|x_i\|^4=(\operatorname{tr}\Sigma)^2+2\operatorname{tr}\Sigma^2+\kappa_4\sum_j\Sigma_{jj}^2,\qquad \mathbb EG_{ij}^2=\operatorname{tr}\Sigma^2\ \ (i\neq j).
\]
Substituting into the multinomial formula and taking the expectation over the particle matrix (weights deterministic or independent):
\[
\mathbb E\big[\operatorname{tr}S_{r,N}^2-\operatorname{tr}S_{w,N}^2\,\big|\,w\big]=\frac{1-s_{2,N}}{N}\Big[(\operatorname{tr}\Sigma_N)^2+\operatorname{tr}\Sigma_N^2+\kappa_4\sum_j(\Sigma_N)_{jj}^2\Big],
\]
exactly. For the asymptotics, $\frac1d\operatorname{tr}\Sigma_N^2\le\|\Sigma_N\|_{\rm op}\frac{\operatorname{tr}\Sigma_N}{d}=O(1)$ (equivalently $\le\|\Sigma_N\|_{\rm op}^2$) and $\frac1d\sum_j(\Sigma_{N,jj})^2\le\|\Sigma_N\|_{\rm op}^2=O(1)$, both under (A5) at finite $N$ (no limit quantities); hence
\[
\frac1d\mathbb E[\mathrm{cost}\mid w]=c_N\Big(\frac{\operatorname{tr}\Sigma_N}{d}\Big)^2(1-s_{2,N})+O(N^{-1})=c\,a_1^2(1-s_{2,N})+o(1),
\]
where no convergence rate is assumed for $c_N$ or the population first moment, so the total error is only $o(1)$. $\square$

\textbf{Corollary~\ref{main-coro:frob-order} (finite-sample Frobenius ordering).}
Write $\lambda=k+u$ with $L=\sum_ik_i$, $R=\sum_iu_i$, $L,R>0$, and note
$L+R=N$. In these coordinates
$\operatorname{Cov}_{\rm multi}=\operatorname{diag}(\lambda)-\frac1N\lambda\lambda^\top$
and
$\operatorname{Cov}_{\rm res}=\operatorname{diag}(u)-\frac1R uu^\top$,
so, using $\operatorname{diag}(\lambda)-\operatorname{diag}(u)=\operatorname{diag}(k)$,
\[
\operatorname{Cov}_{\rm multi}-\operatorname{Cov}_{\rm res}
=\operatorname{diag}(k)-\tfrac1N\lambda\lambda^\top+\tfrac1R uu^\top.
\]
Expanding the claimed rank-one term with $LR/N$ and using $N=L+R$,
\[
\tfrac{LR}{N}\big(\tfrac kL-\tfrac uR\big)\big(\tfrac kL-\tfrac uR\big)^\top
=\tfrac{R}{LN}kk^\top-\tfrac1N\big(ku^\top+uk^\top\big)+\tfrac{L}{RN}uu^\top,
\]
and the $kk^\top$, $ku^\top+uk^\top$, $uu^\top$ coefficients of
$\operatorname{diag}(k)-\frac1N\lambda\lambda^\top+\frac1Ruu^\top$
minus $\operatorname{diag}(k)$ are respectively $-\frac1N$, $-\frac1N$,
$-\frac1N+\frac1R$, while those of
$-\frac1L kk^\top+\frac{LR}{N}(\cdots)$ are
$-\frac1L+\frac{R}{LN}=-\frac1N$, $-\frac1N$, $\frac{L}{RN}=\frac1R-\frac1N$;
the two sides agree, giving the stated identity. The first term
$\operatorname{diag}(k)-kk^\top/L$ is the covariance of a
$\mathrm{Multinomial}(1;k/L)$ draw, hence positive semidefinite, and the
second term is an explicit positive semidefinite rank-one matrix, so the
difference is positive semidefinite. The boundary cases are immediate: if
$R=0$ then $\operatorname{Cov}_{\rm res}=0$; if $L=0$ then $N=R$ and the two
covariances coincide. The ordering of the conditional expected error energies
follows from the quadratic form of Proposition~\ref{main-prop:cost}, since
$G\circ G\succeq0$ (Schur product theorem) and
$A\succeq B\succeq0$ implies $\operatorname{tr}[MA]\ge\operatorname{tr}[MB]$
for $M\succeq0$. $\square$

%==============================================================================
\section{Proof of Theorem~\ref{main-thm:multi-count} (Multinomial count law)}\label{sec:S4}

\textbf{Bounded test functions.} Represent the counts by independent ancestor indices $J_1,\dots,J_N$ with $\mathbb P(J_a=i)=w_i$. Step 1 (expectation): the total-variation bound for the Poisson approximation is
\[
d_{\rm TV}\big(\operatorname{Bin}(N,\lambda/N),\ \operatorname{Poi}(\lambda)\big)\ \le\ \frac{\lambda^2}{N},
\]
which is $O(M^2/N)$ uniformly on $\lambda\in[0,M]$; the tail above $M$ is controlled by (A4) via $\rho_N((M,\infty))\le M^{-2}\int\lambda^2\rho_N\to M^{-2}r_2$. Hence $\frac1N\sum_i\mathbb E f(n_i)\to\int\mathbb E f(\operatorname{Poi}(\lambda))\rho(d\lambda)$. Step 2 (concentration): replacing one ancestor label changes two counts, so $\frac1N\sum_if(n_i)$ changes by at most $4\|f\|_\infty/N$ (equivalently $2\operatorname{osc}(f)/N$); bounded differences give variance $O(N^{-1})$. (If weights are random, the computations above are conditional on $w$.)

\textbf{Second moment (factorial-moment form).} Let $B_N:=\frac1N\sum_in_i(n_i-1)$, so $\frac1N\sum_in_i^2=1+B_N$. Conditionally on $w$,
\[
\mathbb E[B_N\mid w]=(N-1)\sum_iw_i^2=\Big(1-\frac1N\Big)\int\lambda^2\rho_N(d\lambda),
\]
and expanding the collision indicators gives
\[
\operatorname{Var}(B_N\mid w)\ \le\ 2\sum_iw_i^2+4N\sum_iw_i^3\ \longrightarrow\ 0,
\]
since $N\sum_iw_i^3\le(\max_iw_i)\,N\sum_iw_i^2\to0$ under (A4) (no third-moment assumption on the limiting weight law is needed). Hence $\frac1N\sum_in_i^2\to1+r_2=b_2$ in probability. $\square$

%==============================================================================
\section{Poissonization coupling}\label{sec:S5}

\textbf{Distributional identity.} For independent $n_i'\sim\operatorname{Poi}(\mu_i)$ with $R:=\sum_i\mu_i$ a positive integer and $T=\sum_in_i'$, for any $k$ with $\sum k_i=R$,
\[
\mathbb P(\boldsymbol n'=k\mid T=R)=\frac{\prod_ie^{-\mu_i}\mu_i^{k_i}/k_i!}{e^{-R}R^R/R!}=\frac{R!}{\prod_ik_i!}\prod_i\Big(\frac{\mu_i}{R}\Big)^{k_i}=\mathrm{Multinomial}(R;\mu/R)\{k\}.
\]
\textbf{Coupling (label form).} Let $T\sim\operatorname{Poi}(R)$ independently of an infinite i.i.d.\ label sequence $J_1,J_2,\dots$ with $\mathbb P(J_a=i)=\mu_i/R$. Define $n_i'$ as the number of occurrences of $i$ among the first $T$ labels, and $n_i$ among the first $R$. Then $(n_i')$ are independent $\operatorname{Poi}(\mu_i)$, $(n_i)\sim\mathrm{Multinomial}(R;\mu/R)$, and $\sum_i|n_i-n_i'|=|T-R|$ exactly (the two count vectors differ only in the last $|T-R|$ labels). The matrix difference is a signed sum of $|T-R|$ rank-one terms, so $\operatorname{rank}(S_N-S_N')\le|T-R|$; Tool S.1 gives $\|F^{S_N}-F^{S_N'}\|_\infty\le|T-R|/d$, which tends to zero in probability provided $\sqrt R/d\to0$ (satisfied here since $R\le N$ and $d/N\to c>0$). This coupling controls weak spectral convergence; it does not by itself control second spectral moments. $\square$

%==============================================================================
\section{Proof of Theorem~\ref{main-thm:lsd-id} (LSD, identity population)}\label{sec:S6}

\textbf{Step 1 (Poissonization).} Coupling + Tool S.1: multinomial and independent-Poisson models share the same LSD limit; assume $D_N=\operatorname{diag}(n_i)$ with independent $\operatorname{Poi}(\lambda_i)$ entries.

\textbf{Step 2 (truncation, in-probability two-limit scheme).} Fix $M$ and set $D_N^{(M)}:=\operatorname{diag}(n_i\wedge M)$. Then
\[
\operatorname{rank}\Big(\frac1N Z(D_N-D_N^{(M)})Z^\top\Big)\le\#\{i:n_i>M\},\qquad \frac1N\#\{i:n_i>M\}\xrightarrow{\mathbb P}\nu((M,\infty)),
\]
so for every $\varepsilon>0$,
\[
\mathbb P\Big(\big\|F^{S_{r,N}}-F^{S_{r,N}^{(M)}}\big\|_\infty>\frac{\nu((M,\infty))}{c}+\varepsilon\Big)\to0.
\]
For each fixed $M$, Step 3 identifies the truncated model's limit as $\mathcal T_c(\nu^{(M)})$; the two-limit argument is then performed in Lévy or bounded-Lipschitz distance (not KS distance: the limit may have atoms), letting $M\to\infty$ with $\nu((M,\infty))\to0$ and continuity of the fixed-point law in the population law.

\textbf{Step 3 (Silverstein convergence; random-diagonal transfer).} With $W_N^{(M)}:=\frac1d(D_N^{(M)})^{1/2}Z^\top Z(D_N^{(M)})^{1/2}\in\mathbb R^{N\times N}$: the empirical law of $D_N^{(M)}$ converges in probability, not per realization. From any subsequence, extract a further subsequence along which this convergence is almost sure; conditional on such a realization, apply the deterministic-population covariance theorem \cite{silverstein1995strong,bai2010spectral}; bounded convergence of the conditional error probabilities then yields convergence in probability along the original subsequence. Hence $\mu_{W_N^{(M)}}\Rightarrow\mu_W^{(M)}=\nu^{(M)}\boxtimes\mathrm{MP}(1/c)$ in probability.

\textbf{Step 4 (exact companion bookkeeping).} Let $Q_N^{(M)}:=\frac1NZD_N^{(M)}Z^\top\in\mathbb R^{d\times d}$ and $B_N^{(M)}:=\frac1N(D_N^{(M)})^{1/2}Z^\top Z(D_N^{(M)})^{1/2}=\frac dNW_N^{(M)}\in\mathbb R^{N\times N}$. With $A=Z(D_N^{(M)})^{1/2}$, $B=(D_N^{(M)})^{1/2}Z^\top$, the products $AB=NQ_N^{(M)}$ and $BA=NB_N^{(M)}$ share nonzero spectra, and with $r=\operatorname{rank}$, $c_N:=d/N$:
\[
\mu_{B_N^{(M)}}=c_N\,\mu_{Q_N^{(M)}}+\big(1-c_N\big)\delta_0\qquad\text{(exact, every }N\text{)};
\]
for $c_N>1$ the negative coefficient is cancelled by the forced zero atom of $\mu_{Q_N^{(M)}}$.

\textbf{Step 5 (assembly).} $\mu_{B_N^{(M)}}\Rightarrow D_c(\mu_W^{(M)})$, so the truncated model converges to $\mathcal T_c(\nu^{(M)})$; Step 2's two-limit scheme yields $\mu_Q=\mathcal T_c(\nu)$. $\square$

\textbf{Equivalence of the two forms of Definition~\ref{main-def:Tc}.} The companion transform of the $W$-model satisfies $\underline m_W(z)=c\,m_Q(cz)$ (from Step 4's scaling) and, by the classical companion form of Silverstein's equation with ratio $y=1/c$,
\[
z=-\frac1{\underline m_W(z)}+\frac1c\int\frac{t}{1+t\,\underline m_W(z)}\,\nu(dt).
\]
Substituting $\underline m_W(z)=c\,m_Q(cz)$, multiplying by $c$, and setting $\zeta=cz$ gives
\[
\zeta=-\frac1{m_Q(\zeta)}+\int\frac{t}{1+ct\,m_Q(\zeta)}\,\nu(dt),
\]
the covariance-side fixed-point equation. $\square$

\textbf{Normalization check.} $\eta=\delta_1$: $B_N=\frac1NZ^\top Z$ has LSD $c\,\mathrm{MP}(c)+(1-c)\delta_0$ by Step 4's counting applied to the standard Wishart; independently, the covariance-side equation with $\nu=\delta_1$ reduces to the MP fixed point $czm^2+(z+c-1)m+1=0$. Equating gives $\mathcal T_c(\delta_1)=\mathrm{MP}(c)$ and $D_c(\mathrm{MP}(1/c))=c\,\mathrm{MP}(c)+(1-c)\delta_0$. (Moment recursion: $m_2=1+cb_2$, and $m_3=1+3cb_2+c^2b_3$ whenever the corresponding moments are finite.) $\square$


%==============================================================================
\section{Proof of Theorem~\ref{main-thm:lsd-gen} (General population spectrum; Gaussian case)}\label{sec:S7}

Order of argument: (i) truncate counts to $\nu^{(M)}$ as in \ref{sec:S6}; the truncated $Q_N^{(M)}$ satisfies $\|Q_N^{(M)}\|_{\rm op}\le M\|\frac1NZZ^\top\|_{\rm op}=O_{\mathbb P}(1)$ (not a deterministic uniform bound over Gaussian realizations); (ii) given $D_N^{(M)}$ (independent of $Z_N$ by exogeneity), Gaussianity makes $Q_N^{(M)}$ conditionally orthogonally invariant, hence asymptotically free from the deterministic $\Sigma_N$ \cite{collins2006integration}; the probability modes are aligned as follows: the count-law convergence is in probability (\ref{sec:S6} Step 3's subsequence bridge), the freeness conclusion is applied conditionally, and the final LSD convergence is in probability; (iii) the truncated model's LSD is $H\boxtimes\mathcal T_c(\nu^{(M)})$; (iv) remove the truncation by Tool S.1 and pass through weak continuity of the free multiplicative convolution to $H\boxtimes\mathcal T_c(\nu)$. The matrices $\Sigma_N^{1/2}Q_N\Sigma_N^{1/2}$ and $Q_N^{1/2}\Sigma_NQ_N^{1/2}$ share their nonzero spectrum; for the bounded truncated model this relation, rather than any moment-sharing statement, is what the proof uses. $\square$

%==============================================================================
\section{Proof of Proposition~\ref{main-prop:atom} (Zero-eigenvalue mass)}\label{sec:S8}

\textbf{Rank loss (finite sample).} $S_{r,N}=\frac1N\sum_{i:n_i>0}n_iz_iz_i^\top$ has at most $\#\{i:n_i>0\}$ nonzero rank-one summands, so $\operatorname{rank}(S_{r,N})\le\min(d,\#\{i:n_i>0\})$; equality a.s.\ for continuously distributed $Z$ (general position). Note that a finite-sample rank proportion alone does not exclude a non-vanishing fraction of nonzero eigenvalues accumulating at zero (a vanishing fraction cannot form an atom, but a non-vanishing fraction tending to zero would matter), which is why the LSD statement is proved through the transform.

\textbf{LSD zero mass (companion-law max-atom rule).} With $q:=\nu(\{0\})$, the max-atom rule for free multiplicative convolution \cite{belinschi-atoms,bercovici1993free} applied to the companion law gives
\[
\mu_B(\{0\})=\max\big\{q,\ (1-c)_+\big\},\qquad \mu_B=D_c\big(\nu\boxtimes\mathrm{MP}(1/c)\big),
\]
since $\mathrm{MP}(1/c)$ has zero mass $(1-c)_+$ (present iff $c<1$). Reading off from $\mu_B=c\mu_Q+(1-c)\delta_0$:
\[
\mu_Q(\{0\})=\frac{\max\{q,(1-c)_+\}-(1-c)}{c}=\max\Big(0,\ 1-\frac{1-q}{c}\Big).
\]
For the separable case, $\mu_S(\{0\})=\max\{H(\{0\}),\mu_Q(\{0\})\}$ is the same rule applied to $H\boxtimes\mu_Q$; the alternative rank route would require $\operatorname{rank}(\Sigma_N)/d\to1-H(\{0\})$, which fails under (A5) alone (e.g.\ $\Sigma_N=N^{-1}I_d$), and is therefore not used. $\square$

%==============================================================================
\section{Proof of Proposition~\ref{main-prop:moments} (Spectral moments; convergence in probability)}\label{sec:S9}

\textbf{Expectation identity.} Let $Y:=\Sigma^{1/2}ZDZ^\top\Sigma^{1/2}$ and $S=Y/N$; write everything as $\mathbb E\operatorname{tr}Y^2/(dN^2)$. Conditional Gaussian Wick evaluation given $D$:
\[
\mathbb E\big[\operatorname{tr}Y^2\,\big|\,D\big]=\big((\operatorname{tr}\Sigma)^2+\operatorname{tr}\Sigma^2\big)\sum_id_i^2+\operatorname{tr}\Sigma^2\Big(\sum_id_i\Big)^2.
\]
For non-Gaussian entries an extra term $\kappa_4\sum_id_i^2\sum_j\Sigma_{jj}^2$ appears, which after normalization by $dN^2$ is $O(N^{-1})$ under (A5) together with $\frac1N\sum_id_i^2=O(1)$ (in probability or in expectation, as recorded with each result).

\textbf{Convergence in probability (auxiliary lemma).} \emph{Under the standing assumptions, $\frac1d\operatorname{tr}S_N\xrightarrow{\mathbb P}a_1$ and $\frac1d\operatorname{tr}S_N^2\xrightarrow{\mathbb P}m_{2,H}+c\,a_1^2\cdot(\text{diag second moment})$.}

\emph{First moment.} Given $D$, $\mathbb E_Z\big[\frac1d\operatorname{tr}S_N\,\big|\,D\big]=\frac{\operatorname{tr}\Sigma_N}{d}\cdot\frac{\sum_id_i}{N}$ and
\[
\operatorname{Var}_Z\Big(\frac1d\operatorname{tr}S_N\,\Big|\,D\Big)=\frac{2\operatorname{tr}\Sigma_N^2}{d^2N^2}\sum_id_i^2\ \longrightarrow\ 0,
\]
since $\sum_id_i=N$, $\frac1N\sum_id_i^2=O(1)$, and $\frac1d\operatorname{tr}\Sigma_N^2=O(1)$. Conditional Chebyshev plus convergence of the conditional mean give the claim for both pre and post matrices.

\emph{Second moment, two ingredients.} (i) \textbf{Truncated model.} With $d_i^{(M)}:=d_i\wedge M$,
\[
\big\|S_N^{(M)}\big\|_{\rm op}\le M\,\|\Sigma_N\|_{\rm op}\,\Big\|\frac1NZZ^\top\Big\|_{\rm op}=O_{\mathbb P}(1),
\]
so the truncated model has asymptotically bounded spectrum, and its LSD convergence (to $H\boxtimes\mathcal T_c(\nu^{(M)})$, by \ref{sec:S7}) upgrades to convergence in probability of its second spectral moment, whose limit is
\[
m_{2,H}\Big(\int t\,\nu^{(M)}(dt)\Big)^2+c\,a_1^2\int t^2\,\nu^{(M)}(dt)
\]
(the truncated law has mean $\int t\,\nu^{(M)}(dt)$, generally less than one; the mean-one formula does not apply).

(ii) \textbf{Truncation-error control.} Set $X_{N,M}:=\frac1d\big[\operatorname{tr}S(D)^2-\operatorname{tr}S(D^{(M)})^2\big]$. Since $D\succeq D^{(M)}\succeq0$, one has $S(D)\succeq S(D^{(M)})\succeq0$, and
\[
\operatorname{tr}A^2-\operatorname{tr}B^2=2\operatorname{tr}\big[B(A-B)\big]+\operatorname{tr}(A-B)^2\ge0\qquad(A\succeq B\succeq0),
\]
so $X_{N,M}\ge0$. With $\bar d:=\frac1N\sum_id_i$ and $\bar d_M:=\frac1N\sum_id_i^{(M)}$, the Wick identity applied to both matrices gives
\[
\mathbb E_Z\big[X_{N,M}\,\big|\,D\big]=\frac{(\operatorname{tr}\Sigma)^2+\operatorname{tr}\Sigma^2}{dN}\cdot\frac1N\sum_i\big(d_i^2-(d_i^{(M)})^2\big)+\frac{\operatorname{tr}\Sigma^2}{d}\big(\bar d^2-\bar d_M^2\big),
\]
and both terms vanish as $M\to\infty$ by the second tail moment of the count law (finite second moment; the mean difference $\bar d-\bar d_M\to0$ likewise). To remove the truncation without assuming unconditional integrability of the random tail, let $A_{N,M}(D):=\mathbb E_Z[X_{N,M}\mid D]$; since $X_{N,M}\ge0$, for any $\varepsilon,\delta>0$,
\[
\mathbb P(X_{N,M}>\varepsilon)\ \le\ \mathbb P(A_{N,M}>\delta)+\frac{\delta}{\varepsilon},
\]
because $\mathbb P(X>\varepsilon)\le\mathbb P(A>\delta)+\mathbb P(X>\varepsilon,\,A\le\delta)$ and the second term is at most $\mathbb E[A\mathbf 1_{A\le\delta}]/\varepsilon\le\delta/\varepsilon$. First let $N\to\infty$ at fixed $M$, then $M\to\infty$, then $\delta\downarrow0$; the truncation is removed.

The limit value follows from the expectation identity: pre-resampling ($D=\Lambda_N$): diagonal second moment $\to r_2$; post-resampling ($D=D_N$): $\to b_2=1+r_2$ (Theorem~\ref{main-thm:multi-count}). Increment $c\,a_1^2$. $\square$

%==============================================================================
\section{Proof of Proposition~\ref{main-prop:ess} (ESS; occupancy bound; the pair)}\label{sec:S10}

\textbf{ESS.} $N/\mathrm{ESS}=N\sum_iw_i^2\to r_2$, equivalently $\mathrm{ESS}/N\to1/r_2$; given $c$ and $H$, it determines the multinomial limiting second spectral moment (Proposition~\ref{main-prop:moments}), but not the full limiting spectral distribution.

\textbf{Occupancy bound (one line).} For $K\sim\MixPoi(\rho)$, $\mathbb EK=1$ and $\mathbb EK^2=1+r_2$; by Cauchy--Schwarz,
\[
1=(\mathbb EK)^2=\big(\mathbb E[K\mathbf 1_{\{K>0\}}]\big)^2\le\mathbb EK^2\cdot\mathbb P(K>0),
\]
so $p_{\rm multi}\ge1/(1+r_2)$. If $r_2=1$ with mean $1$, then $\mathrm{Var}(K)=0$, forcing $\rho=\delta_1$ and $\nu=\operatorname{Poi}(1)$. $\square$

\textbf{The pair.} $\rho^{(A)}=\frac12\delta_{1/4}+\frac12\delta_{7/4}$: mean $1$; $\mathbb EL^2=25/16$; count-law zero mass $\frac12e^{-1/4}+\frac12e^{-7/4}=0.476287$; count third moment $b_3=2.6875+3(25/16)+1=8.375$. $\rho^{(B)}=\frac9{10}\delta_{3/4}+\frac1{10}\delta_{13/4}$: mean $1$; $\mathbb EL^2=25/16$; zero mass $\frac9{10}e^{-3/4}+\frac1{10}e^{-13/4}=0.429007$; $b_3=3.8125+3(25/16)+1=9.5$. Identical ESS $=16N/25$ exactly at multiples of $10$. Since $m_3(\mathcal T_c(\nu))=1+3cb_2+c^2b_3$ (moments finite here) with equal $b_2=1+r_2=41/16$ and different $b_3$,
\[
m_3^{(B)}-m_3^{(A)}=c^2\,(9.5-8.375)=\frac98c^2>0,
\]
so the two limiting spectra differ for every $c>0$—this example is self-contained and does not rely on the general injectivity lemma (\ref{sec:S21}). At $c=1$, $\mu_Q(\{0\})=\nu(\{0\})$ witnesses $0.476287$ vs.\ $0.429007$. (Differences in third spectral moments establish different distributions; they do not by themselves establish different support edges.) $\square$

%==============================================================================
\section{Proof of Corollary~\ref{main-coro:deff} (Effective dimension; IPR counterexample)}\label{sec:S11}

Under the assumptions of Proposition~\ref{main-prop:moments}, with $a_1>0$: \ref{sec:S9}'s convergence in probability of the normalized first and second moments, plus continuity of $\deff/d$ at the positive limit point, gives the two displayed limits and the explicit gap. No finite-sample ordering is claimed, because of the random denominator: with $d=N=2$, $x_1=3e_1$, $x_2=e_2$, $w=(0.9,0.1)$, $\operatorname{IPR}(S_w)=(8.1^2+0.1^2)/8.2^2\approx0.97591$, while multinomial resampling gives counts $(2,0),(1,1),(0,2)$ with probabilities $0.81,0.18,0.01$ and IPR values $1,0.82,1$, hence the \textbf{conditional expectation} $\mathbb E[\operatorname{IPR}(S_{r,N})\mid X,w]=0.82+0.18\times0.82=0.9676$—strictly smaller. $\square$

%==============================================================================
\section{Proof of Theorem~\ref{main-thm:rankgap} (Rank-and-spectral-gap transition; critical touching)}\label{sec:S12}

\textbf{Active-column reduction and domination.} With $I_N:=\{i:n_i>0\}$,
\[
S_{r,N}=\frac1N\sum_{i\in I_N}n_i\,z_iz_i^\top=\frac1N Z_{I_N}D_{\rm act}Z_{I_N}^\top,\qquad D_{\rm act}\succeq I_{M_N},
\]
since active counts are at least $1$; hence $S_{r,N}\succeq\widetilde W_N:=\frac1NZ_{I_N}Z_{I_N}^\top$ and Weyl monotonicity applies. Conditional on the active index set, $Z_{I_N}$ has i.i.d.\ entries with the original entry distribution, by independence of the counts and the particle matrix (exogeneity). $M_N/N\to p$ in probability by Theorem~\ref{main-thm:multi-count}. Random $M_N$ is handled by conditioning and a subsequence argument (as in \ref{sec:S6} Step 3).

\textbf{Subcritical $c<p$.} $\widetilde W_N=\frac{M_N}{N}\cdot\frac1{M_N}Z_{I_N}Z_{I_N}^\top$; by Tool S.3, $\lambda_{\min}\big(\frac1{M_N}Z_{I_N}Z_{I_N}^\top\big)\to(1-\sqrt{c/p})^2$ in probability, hence $\lambda_{\min}(\widetilde W_N)\to\frac{M_N}{N}(1-\sqrt{c/p})^2\to(\sqrt p-\sqrt c)^2$ in probability, and Weyl monotonicity yields, for every $\varepsilon>0$, $\mathbb P\big(\lambda_{\min}(S_{r,N})\ge(\sqrt p-\sqrt c)^2-\varepsilon\big)\to1$. Since the limit is positive, $\mathbb P(\operatorname{rank}(S_{r,N})=d)\to1$ follows.

\textbf{Supercritical $c>p$.} Pass to the Gram matrix $\Gamma_N:=\frac1NZ_{I_N}^\top Z_{I_N}=\frac dN\cdot\frac1dZ_{I_N}^\top Z_{I_N}$ ($M_N\times M_N$). The matrix $\frac1dZ_{I_N}^\top Z_{I_N}$ has population dimension $M_N$, sample size $d$, and ratio $M_N/d\to p/c<1$ (note: $c>p$), so by Tool S.3 with $y=p/c<1$, $\lambda_{\min}\big(\frac1dZ_{I_N}^\top Z_{I_N}\big)\to(1-\sqrt{p/c})^2$ in probability; scaling by $d/N\to c$ gives $\lambda_{\min}(\Gamma_N)\to c(1-\sqrt{p/c})^2=(\sqrt c-\sqrt p)^2$ in probability. The Gram lower-edge convergence also implies $\mathbb P(\operatorname{rank}(Z_{I_N})=M_N)\to1$ (a positive limiting smallest eigenvalue forces full column rank with high probability; almost-sure rank equalities require continuous entries). On the event $\operatorname{rank}(Z_{I_N})=M_N$, both $\widetilde W_N$ and $S_{r,N}$ have rank exactly $M_N$ (the latter by domination: $\operatorname{range}(\widetilde W_N)\subseteq\operatorname{range}(S_{r,N})$, and the active expansion gives the reverse inequality), and with the convention $\lambda_1\ge\cdots\ge\lambda_d$ (descending), the smallest positive eigenvalue is $\lambda_{M_N}$; Weyl monotonicity at that index gives
\[
\lambda_{\min}^{+}(S_{r,N})=\lambda_{M_N}(S_{r,N})\ \ge\ \lambda_{M_N}(\widetilde W_N)=\lambda_{\min}(\Gamma_N),
\]
hence $\mathbb P\big(\lambda_{\min}^+(S_{r,N})\ge(\sqrt c-\sqrt p)^2-\varepsilon\big)\to1$. The zero-atom mass is Proposition~\ref{main-prop:atom}.

\textbf{Criticality $c=p$.} The fixed-point identity of Definition~\ref{main-def:Tc}, defined initially on $\mathbb C^+$, extends to the negative real axis by continuity (no singularity for nonnegative spectral measures). Write $p=\nu((0,\infty))$, set $c=p$, take $z=-s<0$, and write $m=m(-s)>0$. The equation rearranges to
\[
sm(-s)=\frac1c\int_{(0,\infty)}\frac{\nu(dt)}{1+ctm(-s)}.
\]
\emph{Step 1 ($m(-s)\to\infty$).} Suppose along some sequence $s_j\downarrow0$ one had $m(-s_j)\le C$. Then $\frac{1}{1+ctm(-s_j)}\ge\frac{1}{1+ctC}$ pointwise, so the right side is bounded below by $\frac1c\int_{(0,\infty)}\frac{\nu(dt)}{1+ctC}>0$, while the left side $s_jm(-s_j)\to0$—contradiction. \emph{Step 2 (no atom).} Now that $m(-s)\to\infty$, the integrand tends to $0$ pointwise and is bounded by $1$; dominated convergence gives $sm(-s)\to0$, and since $sm(-s)=\int\frac{s}{x+s}\mu_Q(dx)\to\mu_Q(\{0\})$, this proves $\mu_Q(\{0\})=0$. \emph{Step 3 (touching).} If the nonzero spectrum were separated from zero by a gap $\varepsilon>0$, all mass would lie in $[\varepsilon,\infty)$ and $m(-s)=\int\frac{\mu_Q(dx)}{x+s}\le1/\varepsilon$, contradicting Step 1. Hence $\inf\operatorname{supp}\mu_Q=0$.

\textbf{Finite-sample statement at criticality.} For every $\varepsilon>0$, $\mu_Q((0,\varepsilon))>0$ by Step 3; choose $f\in C_c((0,\varepsilon))$, $f\ge0$, $\int f\,d\mu_Q>0$. Empirical spectral convergence gives $\int f\,d\mu_N\xrightarrow{\mathbb P}\int f\,d\mu_Q>0$, so with probability tending to one at least one eigenvalue lies in $(0,\varepsilon)$: $\mathbb P(\lambda_{\min}^+(S_{r,N})<\varepsilon)\to1$; nonnegativity gives the other direction, hence $\lambda_{\min}^+(S_{r,N})\xrightarrow{\mathbb P}0$. (No no-eigenvalues-outside-the-support input is needed for this statement; such inputs serve stronger claims, e.g.\ convergence to the true positive lower edge off criticality, or edge rates.)

\textbf{Divergence rate of the Stieltjes transform.} Multiplying the critical identity by $m(-s)$:
\[
s\,m(-s)^2=\frac1c\int_{(0,\infty)}\frac{m(-s)}{1+ctm(-s)}\,\nu(dt).
\]
As $s\downarrow0$, $\frac{m(-s)}{1+ctm(-s)}\to\frac1{ct}$ pointwise and $\frac{m(-s)}{1+ctm(-s)}\le\frac1{ct}$; for integer-valued counts, $t\ge1$ on the support, so $t^{-1}$ is integrable, and dominated convergence gives $s\,m(-s)^2\to\frac1{c^2}\int_{(0,\infty)}\frac1t\,\nu(dt)$, i.e.\ $m(-s)\sim\frac1c\big(\int t^{-1}\nu(dt)\big)^{1/2}s^{-1/2}$. (This is an asymptotic for the Stieltjes transform along the negative real axis; a density-level $x^{-1/2}$ asymptotic would require additional argument.)

\emph{Benchmark scope.} The domination argument discards the additional weights $D_{\rm act}-I$, so it need not identify the exact lower edge. The exact regular lower edge depends on the full count law and is characterized by Proposition~\ref{main-prop:edges}. $\square$

%==============================================================================
\section{Proof of Proposition~\ref{main-prop:edges} (Edges; unbounded support via positive fractions)}\label{sec:S13}

\textbf{Edge parametrization.} By \cite{silverstein1995analysis} (see also \cite[Ch.~6]{bai2010spectral}), the law $\mu_W$ characterized by the Silverstein equation admits the real parametrization $z(u)=-\frac1u+y\int\frac{t}{1+tu}\nu(dt)$, $y=1/c$, on real $u\neq0$ with $-1/u\notin\operatorname{supp}\nu$; the support-complement branches are those with $z'(u)>0$, and regular endpoints arise at boundary critical points of these branches. Since $S_{r,N}$ and $B_N=\frac dNW_N$ share nonzero spectra and $B_N$ is $W_N$ scaled by $d/N\to c$, the regular edges of the nonzero part of $\mu_Q$ are $c\,z(u^*)$. Consistency: $\nu=\delta_1$ gives $c(1\mp\sqrt{1/c})^2=(1\mp\sqrt c)^2$. The critical zero endpoint corresponds to $u\to\infty$ rather than a finite critical point and is treated separately by Theorem~\ref{main-thm:rankgap}(iii). $\square$

\textbf{Unbounded support.} Since $\rho$ has mean $1$, $\rho((0,\infty))>0$, so $\nu=\MixPoi(\rho)$ satisfies $\nu(\{k\})>0$ for every $k\in\mathbb Z_{\ge0}$; write $p_K:=\nu([K,\infty))>0$, noting $p_K\to0$. Fix a level $L>0$ and choose $K$ with $p_K<c/4$ and $Kc/4>L$. Set $A_K:=\{i:n_i\ge K\}$ and $B_{K,N}:=\frac KN Z_{A_K}Z_{A_K}^\top$. Then $|A_K|/N\to p_K$ and, by Tool S.3 (and its Gram form, which gives full column rank of $Z_{A_K}$ with probability tending to one), the smallest positive eigenvalue of $B_{K,N}$ converges in probability to
\[
K\big(\sqrt c-\sqrt{p_K}\big)^2\ \ge\ K\Big(\sqrt c-\frac{\sqrt c}{2}\Big)^2=\frac{Kc}{4}\ >\ L,
\]
while with probability tending to one its nonzero eigenvalues number $|A_K|$, a fraction $|A_K|/d\to p_K/c>0$ of the dimension. Since $S_{r,N}\succeq B_{K,N}$ (PSD order), Weyl monotonicity places at least $|A_K|$ eigenvalues of $S_{r,N}$ above $a_K:=K(\sqrt c-\sqrt{p_K})^2$ with high probability. Now fix $L<\ell<a_K$ and take bounded continuous $f$ with $f(x)=0$ for $x\le L$, $f(x)=1$ for $x\ge\ell$, $0\le f\le1$. With probability tending to one,
\[
\int f\,d\mu_N\ \ge\ \frac{|A_K|}{d},
\]
so by weak convergence $\int f\,d\mu_Q\ge p_K/c>0$; and since $f\le\mathbf 1_{(L,\infty)}$, $\mu_Q((L,\infty))>0$. As $L$ was arbitrary, $\sup\operatorname{supp}\mu_Q=\infty$; the same argument gives $\lambda_{\max}(S_{r,N})\xrightarrow{\mathbb P}\infty$ in probability. Whether the support can split into several intervals for specific $\rho$ is left open. $\square$

%==============================================================================
\section{Proof of Theorem~\ref{main-thm:res-count} (Residual count law; including concentration and all branches)}\label{sec:S14}

\textbf{Setup.} $n_i=K_i+T_i$ with $K_i=\lfloor\lambda_i\rfloor$ deterministic and, when $R_N=\sum_iU_i>0$, $\boldsymbol T\sim\mathrm{Multinomial}(R_N;U/R_N)$; when $R_N=0$, $\boldsymbol T=0$. By (A4R), $\pi_N\Rightarrow\pi$ on $\mathbb Z_{\ge0}\times[0,1]$ and $R_N/N\to\E_\pi U$.

\textbf{Weak convergence, case $\E_\pi U>0$.} Apply Lemma~\ref{main-lem:poisson} with $(\mu_i)=(U_i)$, $R=R_N$: the residual counts couple to independent $\tilde n_i'\sim\operatorname{Poi}(U_i)$ with rank defect $O(\sqrt{R_N})=o(N)$; the diffuse condition on residual weights is automatic ($\max_iU_i/R_N\to0$). For bounded $f:\mathbb Z_{\ge0}\to\mathbb R$, set $g_f(k,u):=\mathbb E[f(k+\operatorname{Poi}(u))]$, which is bounded and continuous on $\mathbb Z_{\ge0}\times[0,1]$; hence by (A4R), $\frac1N\sum_ig_f(K_i,U_i)\to\int g_f\,d\pi$—this is the precise point at which the joint convergence $\pi_N\Rightarrow\pi$ (and not merely $\rho_N\Rightarrow\rho$) is needed. Concentration follows from the bounded-differences argument on the residual ancestor indices. Hence $\nu_N\Rightarrow\mathcal L(K+P_U)$.

\textbf{Weak convergence, case $\E_\pi U=0$.} Then $R_N/N\to0$, and
\[
\#\{i:n_i\neq K_i\}\ \le\ R_N\ =\ o(N),
\]
so for bounded $f$, $\big|\frac1N\sum_if(n_i)-\frac1N\sum_if(K_i)\big|\le2\|f\|_\infty R_N/N\to0$; by (A4R), $\pi_N\Rightarrow\pi$, and since $\E_\pi U=0$ forces $U=0$ $\pi$-a.s., $\mathcal L(K)=\mathcal L(K+U)=\rho$.

\textbf{Second-moment convergence (all regimes).} The following expectation and variance bounds hold for every $R_N>0$, regardless of whether $R_N/N$ tends to zero. Decompose
\[
\frac1N\sum_in_i^2=\frac1N\sum_iK_i^2+\frac2N\sum_iK_iT_i+\frac1N\sum_iT_i^2.
\]
The exact expectation (for $R_N>0$) is
\[
\frac1N\sum_i\mathbb E n_i^2=\frac1N\sum_i\lambda_i^2+\frac{R_N}{N}-\frac1{NR_N}\sum_iU_i^2,\qquad \Big|\frac1{NR_N}\sum_iU_i^2\Big|\le\frac1N,
\]
since $\mathbb ET_i^2=U_i(1-U_i/R_N)+U_i^2$ and $K_i^2+2K_iU_i+U_i^2=\lambda_i^2$. For the variance: multinomial variances satisfy $\operatorname{Var}(T_i)\le U_i$ and cross-covariances are nonpositive, so with $q_i=U_i/R_N$,
\[
\operatorname{Var}\Big(\frac2N\sum_iK_iT_i\Big)\ \le\ \frac4{N^2}\sum_iU_iK_i^2\ \le\ \frac4{N^2}\sum_i\lambda_i^2\ =\ O(N^{-1}),
\]
and the collision computation of \ref{sec:S4} with $R_N$ ancestor draws and normalization $N$ (using $\sum_iq_i^2\le1/R_N$, $\sum_iq_i^3\le1/R_N^2$) gives
\[
\operatorname{Var}\Big(\frac1N\sum_iT_i^2\Big)\ \le\ C\,\frac{R_N}{N^2}\ \le\ \frac CN.
\]
The two random terms are not independent, but independence is not needed: $\operatorname{Var}(A+B)\le2\operatorname{Var}(A)+2\operatorname{Var}(B)$. Hence $\frac1N\sum_in_i^2\to r_2+\E_\pi U$ in probability in all regimes ($R_N=0$ is deterministic and immediate). $\square$

%==============================================================================
\section{Proof of Theorem~\ref{main-thm:res-lsd} (Residual spectral transform)}\label{sec:S15}

$\nu_{\rm res}$ has mean $\E[K+U]=1$ and finite second moment $r_2+\E_\pi U$. Silverstein's theorem transfers to a random diagonal $D_N$ that is independent of $Z_N$ with converging empirical law—it does not require the diagonal entries to be independent among themselves. Hence one may work with the raw residual count diagonal directly: (i) truncate the \textbf{total} counts, $D_N^{(M)}:=\operatorname{diag}(n_i\wedge M)$ (truncating only the Poisson part would not bound the diagonal, since $K_i$ may be unbounded); (ii) by Theorem~\ref{main-thm:res-count}, the empirical law of $D_N^{(M)}$ converges to $\nu_{\rm res}^{(M)}$; (iii) conditional on $D_N^{(M)}$, apply the covariance-spectrum theorem (with the subsequence bridge of \ref{sec:S6} Step 3); (iv) remove the truncation by the rank bound. This unifies the $\E_\pi U=0$ and $\E_\pi U>0$ cases; Poissonization is retained as a count-law tool only. (Two distinct truncation controls are used and should not be conflated: LSD rank truncation requires only the tail probability $\nu((M,\infty))\to0$; second-tail-moment control is needed for the empirical-moment arguments of \ref{sec:S9}.) The Gaussian separable extension (\ref{sec:S7}) is unchanged. $\square$

%==============================================================================
\section{Proof of Corollary~\ref{main-coro:res-atom} (Residual zero atom and threshold)}\label{sec:S16}

Under $\nu_{\rm res}=\mathcal L(K+P_U)$, a zero count requires $K=0$ and $P_U=0$; on $\{K=0\}$ the latter has probability $e^{-U}$ (with $U=1$ allowed in the limit, contributing $e^{-1}$). Hence $\nu_{\rm res}(\{0\})=\E_\pi[\mathbf 1_{\{K=0\}}e^{-U}]$ and $p_{\rm res}(\pi)=1-\nu_{\rm res}(\{0\})$. Theorem~\ref{main-thm:rankgap}'s transfer: under the same identity-population, independence, and entry assumptions, the noncritical proof transfers through active-column domination, while the critical proof transfers through the residual LSD and its fixed-point equation. (For the occupancy step, $\mathbf 1_{\{k>0\}}$ is continuous on the discrete count space; if the measures are placed on the real half-axis, replace it by a bounded continuous function agreeing with it at integers.) $\square$

%==============================================================================
\section{Proof of Proposition~\ref{main-prop:res-cost} (Residual second-moment cost)}\label{sec:S17}

The count-level identities $\E[K+P_U\mid K,U]=K+U$ and $\Var(K+P_U\mid K,U)=U$ give $b_2(\nu_{\rm res})=\E[(K+U)^2]+\E[U]=r_2+\E_\pi U$. The spectral-moment statement then follows from the residual version of \ref{sec:S9}'s argument (the truncation route of \ref{sec:S9} applies verbatim once the count second-moment convergence of \ref{sec:S14} is available, so the claim is an empirical-moment convergence, not merely a computation of the limit law's second moment). Strictness: $\E_\pi U\le1$ always (equality possible: the endpoint $\pi=\delta_{(0,1)}$), so the inequality is $\le c\,a_1^2$, strict iff $a_1>0$ and $\E_\pi U<1$; for $a_1>0$, equality with the multinomial cost occurs precisely at $\pi=\delta_{(0,1)}$, since $\E_\pi(K+U)=1$ and $\E_\pi U=1$ jointly force $U=1$, $K=0$ $\pi$-a.s. Vanishing: the cost vanishes iff $a_1=0$ or $\E_\pi U=0$; assuming $a_1>0$, iff $\pi(\{U=0\})=1$—a property of the limit law, not of every finite sample (a vanishing fraction of exceptional weights is allowed). Exact finite-$N$ zero cost: whenever counts are deterministic (in particular when every $\lambda_i\in\mathbb Z_{\ge0}$, giving $R_N=0$ and $n_i=\lambda_i$), or under the geometric degeneracy of \ref{sec:S3}. $\square$

%==============================================================================
\section{Proof of Corollary~\ref{main-coro:rank-order} (Rank-preservation ordering, two levels)}\label{sec:S18}

The multinomial zero mass under the joint description: $\nu_{\rm multi}(\{0\})=\E_\pi[e^{-(K+U)}]=\E_\pi[e^{-L}]$. The residual zero mass (\ref{sec:S16}): $\nu_{\rm res}(\{0\})=\E_\pi[\mathbf 1_{\{K=0\}}e^{-U}]$. Difference:
\[
\nu_{\rm multi}(\{0\})-\nu_{\rm res}(\{0\})=\E_\pi\big[\mathbf 1_{\{K\ge1\}}e^{-(K+U)}\big]\ \ge\ 0,
\]
strict iff $\pi(K\ge1)>0$ (the integrand is strictly positive where $\mathbf 1_{\{K\ge1\}}=1$). The spectral ordering $\mu_{\rm res}(\{0\})\le\mu_{\rm multi}(\{0\})$ follows since $q\mapsto\max(0,1-(1-q)/c)$ is nondecreasing. For the spectral strictness, write $q_{\rm multi}:=\big(1-p_{\rm multi}/c\big)_+$ and $h_0:=H(\{0\})$; then $\mu_{\rm res}(\{0\})<\mu_{\rm multi}(\{0\})$ iff $\pi(K\ge1)>0$ and $q_{\rm multi}>h_0$ (in the identity-population case $h_0=0$, reducing to $c>p_{\rm multi}$; for general $H$ the inequality can be masked by a common atom $h_0$). In the window $c\in(p_{\rm multi},p_{\rm res})$, multinomial resampling carries a macroscopic zero atom while residual resampling is of full rank with probability tending to one (identity population). $\square$

%==============================================================================
\section{Proof of Theorem~\ref{main-thm:classification} (Classification of residual limits under $\rho=\delta_1$)}\label{sec:S19}

\textbf{Necessity.} If $\rho_N\Rightarrow\delta_1$, then (A4R) forces $\pi$ to be supported on $\{(k,u):k+u=1\}$, since the push-forward under the continuous map $(k,u)\mapsto k+u$ must be $\delta_1$. Since $k\in\mathbb Z_{\ge0}$ and $u\in[0,1]$, the line $k+u=1$ meets the product set only at $(1,0)$ and $(0,1)$. Hence $\pi=(1-\alpha)\delta_{(1,0)}+\alpha\delta_{(0,1)}$ for some $\alpha\in[0,1]$, and Theorem~\ref{main-thm:res-count} gives $\nu_{\rm res}=(1-\alpha)\delta_1+\alpha\operatorname{Poi}(1)$, with $\E_\pi U=\alpha$, $\nu_{\rm res}(\{0\})=\alpha e^{-1}$, $p_{\rm res}=1-\alpha/e$, and second-moment cost $c\,a_1^2\alpha$ (Proposition~\ref{main-prop:res-cost}). $\square$

\textbf{Realizability.} (i) $\alpha=0$: exactly uniform weights ($K_i=1$, $U_i=0$, $R_N=0$). (ii) $\alpha=1$ (bounded two-level endpoint construction): let $\lambda_i=1-\frac1N$ for $i<N$ and $\lambda_N=2-\frac1N$. Then $\sum_i\lambda_i=N$ exactly, $\max_i\lambda_i<2$, $\frac1N\sum_i\lambda_i^2=1+\frac1N-\frac1{N^2}\to1$, and $\rho_N\Rightarrow\delta_1$; the first $N-1$ particles have $(\lfloor\lambda_i\rfloor,\operatorname{frac}(\lambda_i))=(0,1-\frac1N)\to(0,1)$, and the last has proportion $1/N\to0$, so $\pi_N\Rightarrow\delta_{(0,1)}$; Theorem~\ref{main-thm:res-count} then yields $\nu_{\rm res}=\operatorname{Poi}(1)$ directly, with no informal per-particle argument. (iii) $0<\alpha<1$: let $\alpha_N\to\alpha$ with $\alpha_NN$ integral and $\varepsilon_N\downarrow0$; put a fraction $\alpha_N$ of particles at $1-\varepsilon_N$ and the rest at $1+\frac{\alpha_N}{1-\alpha_N}\varepsilon_N$. The mean is exactly one: $\alpha_N(1-\varepsilon_N)+(1-\alpha_N)\big(1+\frac{\alpha_N}{1-\alpha_N}\varepsilon_N\big)=1$; the exact second moment is
\[
\frac1N\sum_i\lambda_i^2=1+\frac{\alpha_N}{1-\alpha_N}\,\varepsilon_N^2,
\]
and the exact residual fraction is $R_N/N=\alpha_N$ (both used in experiment E4); the joint limit is $(1-\alpha)\delta_{(1,0)}+\alpha\delta_{(0,1)}$. All constructions satisfy both (A4) and (A4R). $\square$

%==============================================================================
\section{Proof of Corollary~\ref{main-coro:uniform-contrast} (Exact-uniform-weight contrast)}\label{sec:S20}

For $w_i=1/N$ exactly: residual resampling leaves the second-moment matrix unchanged at every finite $N$ ($n_i\equiv1$ deterministically); in the identity-population setting its LSD is therefore $\mathrm{MP}(c)$ (Theorem~\ref{main-thm:res-lsd} with $\eta=\delta_1$, normalization check of \ref{sec:S6}), with zero cost and $p_{\rm res}=1$; in the general Gaussian population model the unchanged matrix has LSD $H\boxtimes\mathrm{MP}(c)$. Multinomial resampling gives $\operatorname{Poi}(1)$, the distorted $\mathcal T_c(\operatorname{Poi}(1))$, cost $c\,a_1^2$, and $p_{\rm multi}=1-e^{-1}$. Both sides have $r_2=1$. The interpolation clause is Theorem~\ref{main-thm:classification}. $\square$

%==============================================================================
\section{Lemma (Injectivity of $\mathcal T_c$; self-contained, all nonnegative laws)}\label{sec:S21}

\textbf{Statement.} For every $c>0$, $\mathcal T_c$ is injective on the class of \textbf{all} probability laws on $[0,\infty)$.

\textbf{Proof.} Suppose $\mathcal T_c(\nu)=\mathcal T_c(\nu')$ and let $m$ be the common Stieltjes transform. For every $z$ in the upper half-plane, Definition~\ref{main-def:Tc}'s fixed-point equation gives
\[
\int\frac{t}{1+ctm(z)}\,\nu(dt)\ =\ \int\frac{t}{1+ctm(z)}\,\nu'(dt).
\]
Take $z=-s<0$ and set $r:=cm(-s)>0$. The map $s\mapsto cm(-s)$ is continuous and strictly decreasing on any interval of the negative real axis ($\frac{d}{ds}m(-s)=-\int\frac{\mu(dx)}{(x+s)^2}<0$), so as $s$ varies over a nondegenerate interval, $r$ varies over a nondegenerate interval $I\subset(0,\infty)$—that is all the proof needs. Define
\[
F_\nu(r):=\int\frac{t}{1+rt}\,\nu(dt)\ =\ \frac1r-\frac1{r^2}\,m_\nu\Big(-\frac1r\Big),
\]
the identity holding since $\frac1r\int\frac{\nu(dx)}{1+rx}=\frac1r-\frac1{r^2}m_\nu(-1/r)$ by direct algebra. Equality of $F_\nu$ and $F_{\nu'}$ on $I$ gives equality of $m_\nu$ and $m_{\nu'}$ on the corresponding negative interval; the identity theorem for analytic functions extends the equality to the common domain of analyticity, and uniqueness of the Stieltjes transform gives $\nu=\nu'$. (This route needs no mean-one assumption, no weighted measure, and no separate recovery of the atom at zero; the earlier cancellation-by-$S$-transform argument is retained only as a remark.) $\square$

%==============================================================================
\section{Additional References for the Supplement}\label{sec:S22}

[S1] J.~W.~Silverstein and S.~Choi, ``Analysis of the limiting spectral distribution of large dimensional random matrices,'' \emph{J. Multivariate Anal.}, vol.~54, no.~2, pp.~295--309, 1995 (main-text \cite{silverstein1995analysis}; edge parametrization).
[S2] S.~Belinschi, ``The atoms of the free multiplicative convolution of two probability distributions,'' \emph{Integral Equations and Operator Theory}, vol.~46, no.~4, pp.~377--386, 2003 (main-text \cite{belinschi-atoms}; max-atom rule used in S.\ref{sec:S8}).
[S3] Z.~D.~Bai and Y.~Q.~Yin, ``Limit of the smallest eigenvalue of a large dimensional sample covariance matrix,'' \emph{Ann. Probab.}, vol.~21, no.~3, pp.~1275--1294, 1993 (main-text \cite{bai1993limit}; Tool S.3).
[S4] H.~Bercovici and D.~Voiculescu, ``Free convolution of measures with unbounded support,'' \emph{Indiana Univ. Math. J.}, vol.~42, no.~3, pp.~733--773, 1993 (main-text \cite{bercovici1993free}).
